\documentclass[a4paper,12pt]{report}
\usepackage{styles/fbe_tez}
\usepackage[utf8x]{inputenc} % To use Unicode (e.g. Turkish) characters
\renewcommand{\labelenumi}{(\roman{enumi})}
\usepackage{amsmath, amsthm, amssymb}
 % Some extra symbols
\usepackage[bottom]{footmisc}
\usepackage{cite}
\usepackage{url}
\usepackage{graphicx}
\usepackage{longtable}
\graphicspath{{figures/}} % Graphics will be here

\usepackage{multirow}
\usepackage{subfigure}
\usepackage{algorithm}
\usepackage{algorithmic}

\begin{document}

% Title Page
\title{CMPE 492 \\ Stablecoin Flow Analysis on Blockchains}
\author{
Celil Ozkan \\ \\
Advisor: \\ 
Prof. Can Ozturan
}
\date{}
\maketitle{}
\pagenumbering{roman}
\tableofcontents

\chapter{INTRODUCTION}
\pagenumbering{arabic}

\section{Broad Impact}
Stablecoins have become a critical infrastructure component in digital finance. They are used for
cross-border settlement, centralized exchange (CEX) operations, decentralized exchange (DEX)
liquidity, lending protocols, and cross-chain bridge transfers. Because stablecoin activity spans
multiple blockchain networks, understanding where value flows and which addresses act as hubs is
important for both technical and economic analysis.

This project focuses on extracting ERC-20 transfer events for selected stablecoins from
EVM-compatible networks and building a data foundation for graph-based flow analysis. The
repository already implements the data collection layer (RPC fetcher + storage) and prepares the
next stages (address classification and graph construction). Even this initial stage has clear impact:
it enables reproducible, large-scale transaction flow studies without depending on closed datasets.

Potential impact areas are:
\begin{itemize}
\item Measuring concentration of liquidity and systemic dependency on specific protocols.
\item Identifying major transfer corridors between users, exchanges, bridges, and protocol contracts.
\item Supporting transparent research on cross-network asset movement.
\item Providing a practical educational pipeline for blockchain data engineering.
\end{itemize}

\section{Ethical Considerations}
This work only uses publicly available on-chain data (event logs and metadata) and does not attempt
to deanonymize real-world identities. However, ethical risks still exist:
\begin{itemize}
\item \textbf{Misclassification risk:} Address labels (e.g., CEX, DEX, bridge) may be incomplete or
outdated, leading to incorrect interpretation.
\item \textbf{Attribution risk:} A wallet may represent multiple actors, so behavioral claims should be
probabilistic rather than absolute.
\item \textbf{Privacy expectations:} Although blockchain data is public, analysis should avoid publishing
harmful, person-targeting conclusions.
\item \textbf{Method transparency:} All assumptions, filtering choices, and known limitations must be
clearly documented.
\end{itemize}

To reduce these risks, the project keeps raw transfer evidence, stores exact chain/log identifiers,
and emphasizes reproducibility and careful interpretation.

\chapter{PROJECT DEFINITION AND PLANNING}
\section{Project Definition}
The objective is to build a modular pipeline that collects stablecoin transfer events and prepares them
for network analysis across multiple EVM chains.

\textbf{Scope in repository:}
\begin{itemize}
\item Networks: Ethereum, Polygon, Arbitrum, Avalanche, Optimism.
\item Tokens: USDT, USDC, DAI (all selected chains), XAUT (Ethereum).
\item Data source: Infura JSON-RPC methods (mainly \texttt{eth\_getLogs} and \texttt{eth\_blockNumber}).
\item Storage: PostgreSQL with partitioned transfer table and progress tracking table.
\item Interface: CLI command for automatic/manual fetching.
\end{itemize}

The long-term project goal is to build user-to-user, user-to-protocol, and cross-chain flow graphs,
with address-type classification as an intermediate stage.

\section{Project Planning}
The implementation plan is organized in three phases:
\begin{enumerate}
\item \textbf{Fetch Phase (implemented):} Build and validate robust event extraction and database
storage.
\item \textbf{Analysis Phase (ongoing):} Add address classification and graph construction.
\item \textbf{Reporting Phase:} Summarize findings, limitations, and future extensions.
\end{enumerate}

\subsection{Project Time and Resource Estimation}
Within the semester timeline, the main effort was split as follows:
\begin{itemize}
\item 30\%: Architecture decisions (network/token selection, schema design, data source migration).
\item 45\%: Core implementation (fetcher, DB writes, partition/index setup, CLI command flow).
\item 15\%: Validation and debugging (Infura 10k limit handling, duplicate handling, progress logs).
\item 10\%: Documentation and report preparation.
\end{itemize}

Main resources:
\begin{itemize}
\item Python toolchain: \texttt{requests}, \texttt{click}, \texttt{psycopg2}, \texttt{pyyaml}, \texttt{dotenv}.
\item PostgreSQL 16 running in Docker.
\item Infura RPC endpoints with API key/secret authentication.
\end{itemize}

\subsection{Success Criteria}
The project defines practical success criteria for the current stage:
\begin{itemize}
\item Correctly fetch and decode ERC-20 \texttt{Transfer} logs for selected tokens.
\item Support both automatic and manual block-range execution via CLI.
\item Store events in relational schema with conflict-safe upsert behavior.
\item Preserve resumability/provenance through chunk-level fetch progress records.
\item Keep architecture modular so classification and graph stages can be added without
rewriting fetch/storage layers.
\end{itemize}

\subsection{Risk Analysis}
Key technical and project risks are:
\begin{itemize}
\item \textbf{RPC result cap (10k logs):} Handled by adaptive chunk resizing when provider returns range
overflow errors.
\item \textbf{Rate limits / provider dependency:} Reduced by shifting from explorer APIs to Infura RPC and
using chunked requests.
\item \textbf{Data duplication:} Mitigated by unique constraint \texttt{(tx\_hash, log\_index, network)} and
\texttt{ON CONFLICT DO NOTHING}.
\item \textbf{Cross-chain comparability:} Controlled by restricting scope to EVM-compatible networks.
\item \textbf{Schedule pressure:} Core pipeline prioritized before optional dashboard and advanced
classification features.
\end{itemize}

\subsection{Team Work (if applicable)}
The repository and implementation are currently maintained by a single student developer.
Coordination was done through regular advisor meetings and weekly technical checkpoints.


\chapter{RELATED WORK}
This project is implementation-oriented and grounded in publicly available blockchain tooling.
Three related areas are particularly relevant:

\textbf{(1) Blockchain data providers and explorers.}
Explorer APIs (Etherscan/Arbiscan/Polygonscan/Snowtrace/Optimistic Etherscan) provide convenient
access but can become restrictive for high-volume extraction due to rate limits and endpoint
limitations. JSON-RPC endpoints (e.g., Infura) provide lower-level access and are better aligned with
event-centric pipelines. Practical RPC method design and usage patterns are discussed in
developer-focused sources such as Sharma's guide \cite{sharma2020} and the official Ethereum
JSON-RPC documentation \cite{ethereum_jsonrpc}.

\textbf{(2) ERC-20 event-based analysis.}
Many token analytics workflows reconstruct transfer networks by parsing \texttt{Transfer} events rather
than tracing full transaction internals. The ERC-20 standard itself is defined in the original Ethereum
Improvement Proposal by Vogelsteller and Buterin \cite{erc20}. This project follows that standard and
stores normalized event records for later graph modeling.

Empirical research also demonstrates that ERC-20 transaction networks can be analyzed effectively
with graph methods for structural insights and prediction tasks \cite{somin2020}.

\textbf{(3) Address classification for blockchain intelligence.}
Graph interpretation quality strongly depends on labeling addresses (EOA, exchange hot wallet,
DEX router/pool, bridge contract, etc.). The current codebase prepares this stage by preserving
normalized sender/receiver/token fields and consistent network metadata.

The references list in this report includes the exact technical documentation used during
implementation (RPC method references and chain explorers).

\chapter{METHODOLOGY}
The methodology is a data engineering pipeline designed for reproducibility and extensibility.

\textbf{Step 1 -- Define chain/token configuration.}
Network metadata and token contract addresses are defined in YAML files. This separates
configuration from code and allows controlled scope updates.

\textbf{Step 2 -- Fetch logs by chunked block ranges.}
For each network, the fetcher requests \texttt{eth\_getLogs} over a block interval and selected token
addresses. Different request steps are used per chain to adapt to activity density.

\textbf{Step 3 -- Decode raw logs into normalized transfer records.}
Each log is decoded into a structured record containing block/tx/log identifiers, token fields,
sender/receiver addresses, and decimal-normalized value.

\textbf{Step 4 -- Persist to PostgreSQL incrementally.}
Decoded transfers are written in batches into a partitioned table. Completed block chunks are also
stored in a progress table to support resumable runs and post-hoc auditing.

\textbf{Step 5 -- Prepare analysis-ready representation.}
The stored records are suitable for graph construction where each transfer event corresponds to a
directed weighted edge from \texttt{from\_address} to \texttt{to\_address}. Edge weight is token amount, and
time context is available via block number/timestamp.

Mathematically, for a token-specific transfer graph $G=(V,E)$:
$$
E = \{(u,v,w,t)\,|\,u=\text{from},\;v=\text{to},\;w=\text{value},\;t=\text{block\_time}\}
$$
and aggregate flow between two addresses can be computed as:
$$
W(u,v)=\sum_{e\in E_{u\rightarrow v}} w_e
$$


\chapter{REQUIREMENTS SPECIFICATION}
The requirements are split into functional and non-functional parts.

\section*{Functional Requirements}
\begin{enumerate}
\item System shall fetch ERC-20 \texttt{Transfer} logs for configured token contracts.
\item System shall support multiple EVM networks through a unified interface.
\item System shall provide manual and automatic fetching modes.
\item System shall decode and store transfer fields in a normalized schema.
\item System shall avoid duplicate inserts using deterministic uniqueness constraints.
\item System shall store fetch-progress checkpoints for each processed block chunk.
\item System shall provide query capability for stored transfer data.
\end{enumerate}

\section*{Non-Functional Requirements}
\begin{enumerate}
\item \textbf{Reliability:} interrupted runs should be recoverable from stored progress.
\item \textbf{Scalability:} schema and indexes should support growth across chains and high event volume.
\item \textbf{Maintainability:} configuration files should allow adding/removing chains and tokens with minimal code changes.
\item \textbf{Traceability:} every stored event should preserve source identifiers
(tx hash, log index, block number, network).
\item \textbf{Usability:} command-line execution should be simple for data collection tasks.
\end{enumerate}

\section*{Primary Use Cases}
\begin{itemize}
\item \textbf{UC-1 Fetch Latest Data:} User starts automatic mode and system fetches from last processed
block to latest chain block.
\item \textbf{UC-2 Fetch Custom Interval:} User provides start/end block and system fetches selected range.
\item \textbf{UC-3 Resume After Failure:} User re-runs fetch and system avoids duplicates while continuing
with remaining chunks.
\item \textbf{UC-4 Export for Analysis:} Analyst queries normalized transfer records from PostgreSQL for
classification/graph tasks.
\end{itemize}


\chapter{DESIGN}

\section{Information Structure}
The central relational model consists of two tables:

\textbf{(1) transfers} (partitioned by \texttt{network})
\begin{itemize}
\item Keys and provenance: \texttt{tx\_hash}, \texttt{log\_index}, \texttt{block\_number}, \texttt{block\_hash}
\item Flow fields: \texttt{from\_address}, \texttt{to\_address}, \texttt{value}, \texttt{raw\_value}
\item Asset context: \texttt{token\_symbol}, \texttt{token\_address}
\item Integrity: unique constraint on \texttt{(tx\_hash, log\_index, network)}
\end{itemize}

\textbf{(2) fetch\_progress}
\begin{itemize}
\item Stores processed chunk ranges and log counts per network.
\item Ensures chunk-level uniqueness via \texttt{(network, chunk\_start, chunk\_end)}.
\end{itemize}

Indexing strategy includes block-based, address-based, token-based, and network-based indexes to
support both temporal and entity-centric queries.

\section{Information Flow}
The runtime data flow is:
\begin{enumerate}
\item CLI receives parameters (network + auto/manual + optional start/end).
\item System determines block interval (DB latest block + RPC latest block in auto mode).
\item Fetcher requests \texttt{eth\_getLogs} for selected token addresses and transfer topic.
\item Response logs are decoded and normalized.
\item Batched inserts are written to PostgreSQL with conflict-safe behavior.
\item Processed chunk metadata is written to \texttt{fetch\_progress}.
\item Loop continues until end block is reached.
\end{enumerate}

If the provider reports result-set overflow (more than 10,000 logs), the chunk size is reduced
adaptively and the request is retried.

\section{System Design}
The codebase is modular and currently includes:
\begin{itemize}
\item \texttt{src/cli.py}: top-level command group.
\item \texttt{src/commands/fetch.py}: fetch command argument parsing and mode control.
\item \texttt{src/fetcher\_client.py}: RPC calls, decode logic, adaptive chunk loop.
\item \texttt{src/db.py}: database connection, inserts, reads, progress and utility operations.
\item \texttt{src/helpers.py}: YAML config access, logging, conversion helpers, Infura auth helpers.
\item \texttt{src/enums.py}: chain/token enums and progress data class.
\item \texttt{sql/*.sql}: schema and index creation scripts.
\item \texttt{config/*.yaml}: network and token configuration.
\end{itemize}

This decomposition separates concerns clearly and reduces coupling between command handling,
data retrieval, decode logic, and persistence.

\section{User Interface Design (if applicable)}
The current user interface is command-line based. The CLI provides one implemented command
(\texttt{fetch}) and placeholders for future commands (\texttt{classify}, \texttt{graph}, \texttt{status}).

Current interface benefits:
\begin{itemize}
\item Low operational complexity for data engineers.
\item Easy integration with server-side scheduling (cron, cloud VM, CI jobs).
\item Deterministic invocation for reproducibility.
\end{itemize}

Planned UI extension is a lightweight dashboard to inspect transfer graphs and address categories.

\chapter{IMPLEMENTATION AND TESTING}
\section{Implementation}
The implemented pipeline uses Python with PostgreSQL 16 and Dockerized database runtime.

Important implementation decisions:
\begin{itemize}
\item Switched from explorer APIs to Infura RPC for improved scalability.
\item Added per-network request step values to reduce failure rates under the 10k RPC cap.
\item Implemented adaptive chunk resizing when overflow errors are returned by provider.
\item Inserted data in batches with conflict handling for idempotent re-runs.
\item Partitioned \texttt{transfers} table by network for better query organization and growth handling.
\end{itemize}

The implementation currently covers extraction and storage. Classification and graph generation are
defined as the next stage and partially prepared in code structure.

\section{Testing}
Automated unit tests are not yet included in the repository. Validation has been performed through
manual and integration-style checks:
\begin{itemize}
\item CLI argument validation for auto/manual modes.
\item Correct retrieval of latest block and latest processed DB block in auto mode.
\item Successful decoding of ERC-20 transfer event fields.
\item Duplicate-protection behavior through unique constraints and conflict handling.
\item Correct insertion of chunk metadata into \texttt{fetch\_progress}.
\item Sampling checks on output JSON and database rows for consistency.
\end{itemize}

A key observed behavior is that high-density ranges can hit provider limits; the adaptive chunk
strategy improved successful completion without manual intervention.

\section{Deployment}
Deployment for development is lightweight and reproducible:
\begin{enumerate}
\item Configure environment variables (DB credentials, Infura key/secret).
\item Start PostgreSQL container via Docker Compose.
\item Execute SQL scripts to create schema and indexes.
\item Run CLI fetch command in auto or manual mode.
\end{enumerate}

Example commands:
\begin{itemize}
\item \texttt{docker compose up -d}
\item \texttt{python3 src/cli.py fetch --network ethereum --auto}
\item \texttt{python3 src/cli.py fetch --network polygon --manual --start 50000000 --end 50001000}
\end{itemize}

The architecture is suitable for future cloud deployment (e.g., scheduled fetch jobs on a VM or
container service).

\chapter{RESULTS}
Current repository outputs demonstrate successful event extraction and normalization.

From sample files in the \texttt{output} folder, a total of \textbf{15,151} transfer events were collected.
The observed sample breakdown is:

\begin{table}[h!]
\centering
\caption{Sample extraction summary from repository outputs}
\begin{tabular}{|l|r|l|}
\hline
\textbf{File} & \textbf{Log Count} & \textbf{Block Span} \\\hline
infura\_logs\_20260308\_172243.json & 217 & 24,613,835--24,613,835 \\\hline
infura\_logs\_20260308\_172345.json & 1,394 & 24,613,832--24,613,840 \\\hline
infura\_logs\_20260309\_211817\_batch\_1.json & 10,000 & 50,000,000--50,000,705 \\\hline
infura\_logs\_20260309\_211817\_batch\_2.json & 3,323 & 50,000,705--50,001,000 \\\hline
infura\_logs\_20260309\_211948\_batch\_1.json & 217 & 50,000,000--50,000,010 \\\hline
\end{tabular}
\end{table}

Aggregated token distribution in these samples:
\begin{itemize}
\item USDT: 8,682 events
\item DAI: 5,286 events
\item USDC: 1,167 events
\item XAUT: 16 events
\end{itemize}

Aggregated network distribution in these samples:
\begin{itemize}
\item Ethereum: 1,611 events
\item Polygon: 13,540 events
\end{itemize}

These results verify that the implemented fetch-storage pipeline is functional across different
chains and transfer densities. They also provide a concrete base for the next project stage:
address classification and graph analytics.

\chapter{CONCLUSION}
This project established the core data pipeline for stablecoin flow analysis on EVM-compatible
blockchains. The implemented system fetches ERC-20 transfer logs, decodes them into analysis-ready
records, and stores them in a partitioned PostgreSQL schema with resumable progress tracking.

The current milestone demonstrates reliable extraction and persistence for selected networks and
tokens. This directly enables upcoming research tasks such as address-type classification,
construction of directed weighted flow graphs, and computation of centrality/flow metrics.

Main limitations at this stage are the absence of automated tests, incomplete address labeling, and
no finalized visualization layer. Future work will focus on:
\begin{itemize}
\item Robust address taxonomy (EOA, DEX, CEX, bridge, mint/burn, admin/multisig).
\item Graph generation pipelines and temporal flow analysis.
\item Dashboard/reporting interface for interactive exploration.
\item Expanded evaluation methodology and benchmark scenarios.
\end{itemize}

Overall, the repository already provides a solid and extensible technical foundation for the full
scope of the CMPE 492 stablecoin flow analysis project.

\bibliographystyle{styles/fbe_tez_v11}
\bibliography{references}

\appendix	
\chapter{COMMAND EXAMPLES}
Useful command invocations from the current implementation:

\begin{itemize}
\item \texttt{python3 src/cli.py}
\item \texttt{python3 src/cli.py fetch --network ethereum --auto}
\item \texttt{python3 src/cli.py fetch --network polygon --manual --start 50000000 --end 50001000}
\item \texttt{docker compose up -d}
\end{itemize}

Expected next commands (planned, not yet implemented):
\begin{itemize}
\item \texttt{python3 src/cli.py classify ...}
\item \texttt{python3 src/cli.py graph ...}
\end{itemize}

\end{document}